% =====================================================================
% Sous-chapitre 5.A -- Coeur LangGraph (livrable sous-agent A)
% =====================================================================

\section{Architecture du graphe LangGraph}

Le système de recrutement implémente un superviseur décentralisé via un \texttt{StateGraph} LangGraph : une machine à états typée où l'état global \texttt{GraphState} (voir \texttt{src/state.py}) joue le rôle de tableau noir (\textit{blackboard}) au sens des SMA classiques. Cet état comprend une quinzaine de champs distincts, chacun dédié en écriture à un agent ou à un groupe d'agents : orchestrateur \(\rightarrow\) analyste \(\rightarrow\) pipeline A3 (stratège \(\rightarrow\) collecteur \(\rightarrow\) filtre) \(\rightarrow\) déduplicateur \(\rightarrow\) évaluateurs parallèles \(\rightarrow\) vérificateur \(\rightarrow\) recruteur ou rapport \(\rightarrow\) persistance RAG.

Les nœuds du graphe compilé (\texttt{src/graph.py}) forment une topologie en pipeline avec deux branchements décisionnels significatifs : (i) un \textit{fan-out} vers $N$ évaluateurs parallèles au départ de chaque profil dédupliqué, déclenché par un \texttt{Send()} ; (ii) un \textit{fan-in} par agrégation des scores via un nœud de réduction explicite ; (iii) un routage conditionnel après vérification, qui aiguille soit vers le recruteur, soit vers le rapport en fonction des seuils de score. L'ordre d'exécution respecte une dépendance stricte : huit nœuds séquentiels (A1 à A8) plus les instances parallèles d'A4.

Les champs du \texttt{GraphState} typés \texttt{Annotated[list, operator.add]} (\texttt{profils\_bruts}, \texttt{candidats\_scores}, \texttt{messages\_envoyes}) implémentent des \textit{reducers} qui agrègent les résultats des instances parallèles sans conflit. Le compilateur LangGraph synchronise automatiquement les $N$ branches d'évaluation sur le nœud de réduction, garantissant la cohérence des écritures concurrentes sur l'état partagé.

\section{Patterns de coordination retenus}

Cinq patrons de coordination sont mobilisés dans le graphe, chacun directement transposable à la terminologie SMA classique.

\paragraph{Send/Map-Reduce (fan-out/fan-in).}
La fonction \texttt{route\_to\_evaluateurs()} (voir \texttt{src/graph.py}) crée une liste de $N$ appels \texttt{Send()}, chacun instanciant une copie d'A4 (Évaluateur) avec un candidat distinct et le profil de compétences partagé. Le \textit{reducer} \texttt{operator.add} fusionne les listes \texttt{candidats\_scores} produites par chaque instance. La constante \texttt{MAX\_PROFILS\_PARALLELES} (10 par défaut, dans \texttt{src/config.py}) borne la concurrence afin d'éviter la saturation de l'API LLM locale (Ollama). Le nœud \texttt{reduce\_scores} synchronise la réduction et journalise les statistiques consolidées. La figure~\ref{fig:send_reduce} illustre ce patron.

% Figure : pattern Send / Reduce
\begin{figure}[H]
\centering
\begin{tikzpicture}[node distance=8mm and 8mm, font=\small]
  \node[agent] (a6) {A6 Déduplicateur\\\footnotesize\texttt{profils\_dedupliques}};

  \node[smallagent, below=14mm of a6, xshift=-32mm] (a41) {A4$_1$\\\footnotesize candidat 1};
  \node[smallagent, below=14mm of a6, xshift=-10mm] (a42) {A4$_2$\\\footnotesize candidat 2};
  \node[smallagent, below=14mm of a6, xshift=12mm]  (a43) {A4$_3$\\\footnotesize candidat 3};
  \node[font=\footnotesize, below=14mm of a6, xshift=28mm] (dots) {$\cdots$};
  \node[smallagent, below=14mm of a6, xshift=42mm]  (a4n) {A4$_N$\\\footnotesize candidat $N$};

  \node[agent, below=22mm of a42, text width=4.4cm] (red)
       {Reducer \texttt{operator.add}\\\footnotesize fusion des listes \texttt{candidats\_scores}};

  \node[agent, below=of red] (a5) {A5 Vérificateur};

  \draw[flow] (a6.south) -- node[midway, right=10mm, font=\footnotesize\itshape]{\texttt{Send()} \(\times N\)} (a42.north);
  \draw[flow] (a6.south) -- (a41.north);
  \draw[flow] (a6.south) -- (a43.north);
  \draw[flow] (a6.south) -- (a4n.north);

  \draw[flow] (a41.south) -- (red.north);
  \draw[flow] (a42.south) -- (red.north);
  \draw[flow] (a43.south) -- (red.north);
  \draw[flow] (a4n.south) -- (red.north);

  \draw[flow] (red) -- (a5);

  \node[noteblock, right=18mm of a42, text width=4.6cm] (note)
       {\textbf{Borne :} \texttt{MAX\_PROFILS\_PARALLELES}\\
        Évaluations indépendantes, sans état partagé, fusionnées par le \emph{reducer}.};
\end{tikzpicture}
\caption[Pattern Send / Reduce]{Pattern \emph{Send / Reduce} de LangGraph appliqué à l'évaluation parallèle. Chaque instance d'A4 traite un candidat distinct ; le \emph{reducer} fusionne les listes \texttt{candidats\_scores} sans collision.}
\label{fig:send_reduce}
\end{figure}


\paragraph{Validation pair-à-pair (A4 \(\rightarrow\) A5).}
L'évaluateur (A4) produit des scores bruts dans le champ \texttt{candidats\_scores} ; le vérificateur (A5) reçoit ces scores et croise avec les profils bruts conservés en \texttt{profils\_dedupliques}, puis valide ou invalide chaque candidat indépendamment, en produisant \texttt{candidats\_valides} avec ajustements ou statut « douteux ». Cette séparation formelle des champs impose une vérification transversale obligatoire : A5 contrôle la cohérence intra-CV (dates suspectes, compétences invraisemblables) et inter-candidats (signaux d'alerte, dérives de scoring) sans que le résultat brut d'A4 ne puisse être implicitement réutilisé.

\paragraph{Routage conditionnel par score.}
La fonction \texttt{route\_apres\_verification()} (\texttt{src/graph.py}) applique une logique à deux étages. Si le meilleur candidat dépasse \texttt{SCORE\_SEUIL\_CONTACT} (75 par défaut), le flux est dirigé vers A7 (Recruteur). Sinon, si au moins un candidat dépasse \texttt{SCORE\_SEUIL\_VIABLE} (40), les \texttt{TOP\_N\_RELATIF} meilleurs (3 par défaut) passent en « mode relatif » via A7 avec signalisation explicite. Au-delà, le flux est dirigé directement vers le rapport. Les seuils sont configurables via variables d'environnement.

\paragraph{Interrupt\_before pour le human-in-the-loop.}
Lors de la compilation du graphe, l'option \texttt{interrupt\_before=["recruteur"]} est appliquée au \textit{checkpointer}. Cela suspend l'exécution juste avant le nœud A7, permettant à un humain de valider, d'éditer ou de rejeter la liste de candidats avant tout contact automatisé. La reprise est faite via \texttt{app\_graph.update\_state()} puis nouvel \texttt{invoke(None, config)} : le graphe redémarre exactement à l'endroit où il s'était arrêté.

\paragraph{Mémoire contextuelle (RAG + Persistance).}
Le nœud A8 (\texttt{src/agents/persistance.py}) stocke chaque fiche de poste et ses candidats validés dans ChromaDB, reliés par un \texttt{fiche\_id} stable. L'évaluateur A4 interroge cette mémoire en début de scoring, récupérant les profils similaires évalués pour des postes comparables à titre de calibration inter-runs.

\section{Justification des choix d'implémentation}

\subsection{StateGraph et état typé plutôt que MessageGraph}

\paragraph{Besoin.} Coordonner huit agents spécialisés produisant des artefacts hétérogènes (listes de requêtes, hits bruts, profils dédupliqués, scores détaillés, validations, messages, rapport), avec dépendances complexes et exigence d'analyse \textit{post mortem}.

\paragraph{Alternatives.} Une approche \texttt{MessageGraph} (style chat multi-tour) accumulerait tous les résultats dans une liste \texttt{messages} de LangChain : pratique pour le débogage de prototypes conversationnels, mais sérialisation volumineuse, absence de typage métier, fusion manuelle des résultats parallèles, routage conditionnel maladroit.

\paragraph{Décision.} \texttt{StateGraph} avec \texttt{TypedDict} fortement typé. Chaque champ de \texttt{GraphState} documente qui écrit et qui lit. Les \textit{reducers} \texttt{operator.add} sur les champs parallèles éliminent les collisions de fusion. Les nœuds retournent des dictionnaires partiels, fusionnés automatiquement par le compilateur LangGraph. Le typage statique détecte les fautes de champ dès le développement.

\paragraph{Tradeoffs.} Approche moins flexible qu'un journal de messages plat pour explorer du code dynamique ; exige une architecture bien définie en amont. Compromis assumé : huit agents spécialisés ne forment pas un système conversationnel.

\subsection{Send/Reduce pour l'évaluation parallèle}

\paragraph{Besoin.} Chaque candidat doit être évalué indépendamment (LLM sans état partagé). L'évaluation séquentielle d'un lot de dix profils prendrait de l'ordre de 30 à 40 secondes ; en parallèle, trois à quatre secondes suffisent.

\paragraph{Alternatives.} (i) Boucle interne dans un nœud monolithique ; (ii) orchestration externe à base de threads Python ; (iii) implémentation manuelle d'un \textit{map-reduce} avec fusion ad hoc.

\paragraph{Décision.} Le patron \texttt{Send()} de LangGraph crée implicitement $N$ instances légères du nœud évaluateur, chacune traitant un candidat isolé, avec synchronisation et fusion automatiques. Le \textit{reducer} \texttt{operator.add} agrège les listes sans logique de remplacement.

\paragraph{Tradeoffs.} La concurrence est bornée par \texttt{MAX\_PROFILS\_PARALLELES} pour éviter la surcharge d'Ollama local ; les batches successifs résolvent le cas où plus de candidats doivent être évalués.

\subsection{Interrupt\_before sur A7}

\paragraph{Besoin.} Avant tout envoi de message de contact, un humain doit valider la pertinence de la sélection et le contenu rédigé.

\paragraph{Alternatives.} (i) Pause externe gérée par l'API ; (ii) nœud A7 demandant un \textit{input} synchrone ; (iii) absence de pause, avec audit \textit{post mortem}.

\paragraph{Décision.} L'option \texttt{interrupt\_before} du compilateur LangGraph arrête l'exécution juste avant A7. L'état est entièrement accessible, modifiable et rejouable via le \textit{checkpointer}.

\paragraph{Tradeoffs.} Étape bloquante potentiellement longue ; la transparence et la réduction des erreurs irréversibles justifient le coût dans un contexte RH.

\subsection{Routage conditionnel à seuils absolu et relatif}

\paragraph{Besoin.} Distinguer deux scénarios opérationnels : (i) au moins un candidat clairement au-dessus du seuil (contact automatique) ; (ii) aucun excellent, mais quelques viables (top-N relatif soumis à validation humaine).

\paragraph{Alternatives.} Toujours contacter le top-1 ; contacter tous les candidats au-dessus d'un seuil unique ; deux graphes séparés.

\paragraph{Décision.} Routage dual intégré. \texttt{route\_apres\_verification()} distingue les deux modes et bascule A7 ou rapport en conséquence. Les seuils \texttt{SCORE\_SEUIL\_CONTACT}, \texttt{SCORE\_SEUIL\_VIABLE} et \texttt{TOP\_N\_RELATIF} sont configurables par variables d'environnement.

\paragraph{Tradeoffs.} Logique métier dans le graphe ; justifié par la diversité des fiches de poste et l'asymétrie des marchés selon le métier.

\subsection{Séparation A4 (Évaluateur) / A5 (Vérificateur)}

\paragraph{Besoin.} L'évaluation par LLM est probabiliste (hallucinations, biais de \textit{prompt}, dérives de scoring) ; un contrôle transversal indépendant est nécessaire.

\paragraph{Alternatives.} A4 unique avec auto-vérification ; boucle de \textit{feedback} A4 \(\rightarrow\) A5 \(\rightarrow\) A4 récursive.

\paragraph{Décision.} Deux agents distincts avec champs d'état séparés. A5 reçoit les scores et les profils bruts originaux, et vérifie incohérences et signaux suspects. Patron de \textit{checks and balances} formalisé en SMA classique.

\paragraph{Tradeoffs.} Doublement des appels LLM (latence accrue) ; justifié par la criticité RH d'un mauvais scoring.

\subsection{Ollama local plutôt que API externe (Anthropic, OpenAI)}

\paragraph{Besoin.} Souveraineté des données candidats, coût marginal nul après déploiement, démonstration hors ligne, conformité RGPD.

\paragraph{Alternatives.} Claude API (qualité premium, coût et exposition externe des PII) ; GPT-4 (idem) ; LLaMA local (qualité moindre sur le raisonnement complexe).

\paragraph{Décision.} Ollama via \texttt{init\_chat\_model} de \texttt{langchain\_classic}, avec modèle local configurable (\texttt{OLLAMA\_MODEL=kimi-k2.5:cloud} par défaut). Le fournisseur Ollama est abstrait par LangChain, facilitant un éventuel basculement vers Anthropic ou OpenAI.

\paragraph{Tradeoffs.} Qualité de raisonnement légèrement inférieure aux modèles \textit{frontier}, atténuée par un \textit{prompt engineering} contraint, une température basse (0 pour l'évaluation, 0,7 pour la rédaction) et une structuration JSON stricte des réponses.

\subsection{Découpage en huit agents spécialisés}

\paragraph{Besoin.} Modéliser six phases métier (analyse de fiche, sourcing, déduplication, évaluation, vérification, contact, rapport) plus une phase technique (persistance RAG).

\paragraph{Alternatives.} Monolithe unique ; trois ou quatre agents génériques ; douze sous-agents très spécialisés.

\paragraph{Décision.} Un agent par phase. A1 (superviseur), A2 (analyse), A3a/b/c (sourcing décomposé), A4, A5, A6, A7, A8. Chaque agent lit un sous-ensemble précis du \texttt{GraphState}, écrit ses résultats, et ignore les autres.

\paragraph{Tradeoffs.} Davantage de champs d'état à maintenir ; compense par la testabilité, la traçabilité (logs par agent) et la maintenabilité (modifier un agent ne nécessite pas de refactor systémique).

\subsection{Checkpointer (\texttt{MemorySaver})}

\paragraph{Besoin.} Reprendre l'exécution après interruption HITL ou crash, et exporter l'historique pour audit.

\paragraph{Alternatives.} Absence de \textit{checkpoint} ; base externe (PostgreSQL) ; sérialisation ad hoc.

\paragraph{Décision.} Le compilateur LangGraph reçoit \texttt{checkpointer=MemorySaver()}. Les snapshots d'état sont stockés en mémoire (suffisant pour le démonstrateur). Une variante \texttt{SqliteSaver} permet la persistance entre sessions sans changement de code applicatif.

\paragraph{Tradeoffs.} \texttt{MemorySaver} perd l'historique au redémarrage du processus ; \texttt{SqliteSaver} le persiste mais introduit un fichier de base local.

\section{Extraits commentés}

\begin{lstlisting}[language=Python,caption={Création des instances parallèles d'évaluateur via \texttt{Send()} (\texttt{src/graph.py})}]
def route_to_evaluateurs(state: GraphState) -> list[Send]:
    """Fan-out : envoie chaque profil dédupliqué à une instance d'évaluateur."""
    profils = state.get("profils_dedupliques", [])
    if not profils:
        return [Send("reduce_scores", {})]

    if len(profils) > MAX_PROFILS_PARALLELES:
        profils = profils[:MAX_PROFILS_PARALLELES]

    return [
        Send("evaluateur", {
            "candidat": candidat,
            "profil_competences": state["profil_competences"],
            "fiche_poste": state.get("fiche_poste", ""),
        })
        for candidat in profils
    ]
\end{lstlisting}

Ce patron \texttt{Send()} instancie $N$ nœuds \texttt{evaluateur} isolés, chacun recevant un candidat distinct et les données partagées. LangGraph lance ces $N$ branches en parallèle, attend leur convergence et fusionne les résultats via le \textit{reducer} \texttt{operator.add}.

\begin{lstlisting}[language=Python,caption={Routage conditionnel absolu / relatif (\texttt{src/graph.py})}]
def route_apres_verification(state: GraphState) -> str:
    """Routage après vérification : absolu si meilleur >= seuil, sinon relatif."""
    candidats_valides = state.get("candidats_valides", [])
    if not candidats_valides:
        return "rapport"

    meilleur_score = max(c["score_final"] for c in candidats_valides)

    if meilleur_score >= SCORE_SEUIL_CONTACT:
        return "recruteur"

    viables = [c for c in candidats_valides
               if c["score_final"] >= SCORE_SEUIL_VIABLE]
    if viables:
        return "recruteur"   # Top-N relatif

    return "rapport"
\end{lstlisting}

La logique à deux étages bascule entre contact automatique (mode absolu, seuil strict), top-N relatif (si aucun excellent mais quelques viables) ou rapport sans contact. Les seuils sont configurables.

\begin{lstlisting}[language=Python,caption={Validation pair-à-pair : A5 ré-inspecte avec les profils bruts (\texttt{src/agents/verificateur.py})}]
def verificateur_node(state: GraphState) -> dict:
    """Vérifie les scores de A4 : incohérences dates, CV gonflés, signaux fraude."""
    candidats_scores = state.get("candidats_scores", [])
    llm = init_chat_model(OLLAMA_MODEL, model_provider=OLLAMA_PROVIDER, temperature=0)

    profils_par_id = {p["id"]: p for p in state.get("profils_dedupliques", [])}
    candidats_info = []
    for cs in candidats_scores:
        entry = {
            "candidat_id": cs["candidat_id"],
            "nom": cs["nom"],
            "score_global": cs["score_global"],
            "scores_detail": cs["scores_detail"],
            "resume_evaluateur": cs["resume"],
        }
        profil_brut = profils_par_id.get(cs["candidat_id"])
        if profil_brut:
            entry["profil_brut"] = profil_brut["profil_brut"]
        candidats_info.append(entry)
    # le LLM ré-évalue : cohérence des dates, plausibilité des compétences,
    # détection d'artefacts suspects.
\end{lstlisting}

A5 reçoit les scores produits par A4 et les profils bruts originaux conservés dans un champ distinct du \texttt{GraphState}, ce qui rend possible une vérification transversale des incohérences (dates suspectes, titres vagues, technologies incompatibles). La séparation explicite des champs rend cette validation obligatoire.

\begin{lstlisting}[language=Python,caption={Persistance RAG et lien fiche \(\leftrightarrow\) candidat (\texttt{src/agents/persistance.py})}]
def persistance_node(state: GraphState) -> dict:
    """Stocke fiche + candidats validés en ChromaDB pour RAG futur."""
    candidats_valides = state.get("candidats_valides", [])
    fiche_poste = state.get("fiche_poste", "")

    memoire = get_memoire()

    fiche_id = memoire.ajouter_fiche_poste(fiche_poste)

    profils_idx = {p["id"]: p for p in state.get("profils_dedupliques", [])}

    for candidat in candidats_valides:
        if candidat.get("statut") == "invalide":
            continue   # Ne pas persister le bruit

        profil = profils_idx.get(candidat["candidat_id"], {})
        memoire.ajouter_candidat(
            candidat_id=candidat["candidat_id"],
            nom=candidat["nom"],
            profil_brut=profil.get("profil_brut", "..."),
            score=candidat["score_final"],
            fiche_id=fiche_id,   # lien stable pour requête RAG ultérieure
        )
\end{lstlisting}

A8 constitue le maillon de mémoire : chaque fiche de poste reçoit un identifiant stable (hash du contenu), puis chaque candidat validé est indexé en ChromaDB avec ce \texttt{fiche\_id}. Les exécutions futures peuvent ainsi retrouver les contextes pertinents via \texttt{rechercher\_similaires()} dans A4.
